\chapter*{Project Objective}
\thispagestyle{fancy}
\addcontentsline{toc}{chapter}{Project Objective}
\begingroup
\setlength{\emergencystretch}{2em}

The objective of this project is to establish a traceable analytical and numerical
workflow for normal-incidence reflection from a finite dielectric slab. The study
connects the intrinsic impedance and phase constant of a nonmagnetic dielectric to an
equivalent transmission-line section and then tests the resulting matching condition
using a CST parameter sweep.

The project has five principal technical objectives:

\begin{enumerate}
    \item Derive the input impedance and reflection coefficient of an
    air--dielectric--air slab under normal incidence.

    \item Show that a slab with physical thickness $d=\lambda_0$ becomes matched
    whenever its electrical thickness is an integer multiple of one-half wavelength,
    yielding $\varepsilon_r=(m/2)^2$.

    \item Construct a two-port CST frequency-domain model with electric and magnetic
    side boundaries that approximates TEM-like propagation along the slab axis.

    \item Sweep relative permittivity from 1 to 10 in 0.25 increments and compare the
    locations of the simulated $|S_{11}|$ minima with the analytical prediction.

    \item Examine electric- and magnetic-field evidence at $\varepsilon_r=4$, one of
    the predicted matched cases, while identifying the numerical limitations that
    constrain interpretation of the null depths and port-mode results.
\end{enumerate}

The final outcome is a public engineering report consistent with the other portfolio
projects, together with analytical verification code, machine-readable results,
simulation evidence, and explicit claim boundaries.

\endgroup
